\addtocontents{toc}{\protect\Needspace{3\baselineskip}}
\originalchapter{第三次测验复习：热粒子数与假想新粒子}
以下五组均为 8.286 Fall 2013 Quiz 3 Review Problems 中可公开取得的完整问题。本章沿题目的理想相对论气体、瞬时解耦等假设解答；涉及数值时给出可复核的常数或表达式。题目在教学上用“添加粒子”辨认能量、粒子数、熵自由度三者不同，不能只改一处 $g$。

\src{8.286 Quiz 3 Review Problem 3；官方解答在同一复习册。}
\begin{exercise}
今天 CMB 温度为 $2.7\,\mathrm K$，求光子数密度；再求三味大爆炸遗留热中微子及反中微子的总数密度。采用标准瞬时中微子解耦和电子正电子湮灭近似。
\end{exercise}
\begin{solution}
零化学势的玻色气体有 $n_\gamma=2\zeta(3)(\kb T_\gamma)^3/[\pi^2(\hbar c)^3]$，代 $2.7\,\mathrm K$ 得约 $399\,\mathrm{cm}^{-3}$；若用 $2.725\,\mathrm K$ 数字略变。单个费米态的数密度是同温单个玻色态的 $3/4$；三味各含中微子与反中微子共六态，粒子数自由度 $g_{\nu,n}=6(3/4)=9/2$。熵守恒给 $T_\nu/T_\gamma=(4/11)^{1/3}$，所以
\[
\frac{n_{\nu+\bar\nu}}{n_\gamma}=\frac{(9/2)}2\frac4{11}=\frac9{11},
\qquad n_{\nu+\bar\nu}\simeq326\,\mathrm{cm}^{-3}.
\]
这里把三味的粒子和反粒子都计入，不把 $7/8$ 的\emph{能量}因子误放进数密度。
\end{solution}

\src{8.286 Quiz 3 Review Problem 4(a)--(e), 原题为1998年 Quiz 3 Q4；官方复习解答。}
\begin{exercise}
对温度 $T$ 的黑体：(a) 求每光子平均能量；(b) 求每光子平均熵；(c) 换成一个自旋为零、仅一个自旋态的无质量玻色子，两结果怎样变；(d) 换成仅一个自旋态的无质量电子中微子，求每粒子平均能量；(e) 求其每粒子平均熵。可用 $\zeta(3)$ 保留精确表达式。
\end{exercise}
\begin{solution}
对玻色子，自旋简并度 $g_b$ 的 $u=g_b\pi^2(\kb T)^4/[30(\hbar c)^3]$，$n=g_b\zeta(3)(\kb T)^3/[\pi^2(\hbar c)^3]$，$s=4u/(3T)$。所以
\[
\frac un=\frac{\pi^4}{30\zeta(3)}\kb T\simeq2.701\kb T,\quad
\frac sn=\frac{2\pi^4}{45\zeta(3)}\kb\simeq3.602\kb.
\]
两者都不依赖 $g_b$，故(c)不变。单个费米态的能量、数目因子分别为 $7/8,3/4$，故
\[
\frac{u_\nu}{n_\nu}=\frac{7\pi^4}{180\zeta(3)}\kb T\simeq3.151\kb T,\quad
\frac{s_\nu}{n_\nu}=\frac{7\pi^4}{135\zeta(3)}\kb\simeq4.202\kb .
\]
这里以 $\mu_\nu=0$、质量可忽略为条件；同时计入同温反中微子时，分子分母同倍增，每粒子平均量不变。
\end{solution}

\src{8.286 Quiz 3 Review Problem 5(a)--(e), 原题为1992年 Quiz 3 Q2；官方复习解答。}
\begin{exercise}
假设另有一类与电子性质相同、静止能量 $0.750\,\mathrm{MeV}$ 的轻子及反粒子，题目称“8.286ion”。(a) 在 $\kb T=3\,\mathrm{MeV}$ 时求它们合计每立方米数密度。(b) 在平直辐射主导宇宙 $t=0.01\,\mathrm s$ 求质量密度，以 $\mathrm{g\,cm^{-3}}$ 表示，假设 $0.75\,\mathrm{MeV}\ll\kb T\ll100\,\mathrm{MeV}$，光子、电子对、三味中微子与新轻子对都相对论性。(c) 同时刻求 $\kb T$，以 MeV 表示。(d) 对氦合成方向有何影响？(e) 若中微子在新粒子消失前解耦，求今天 $T_\nu/T_\gamma$。
\end{exercise}
\begin{solution}
(a) 粒子与反粒子各两自旋，共四态，数目因子 $g_{n}=4(3/4)=3$，因此
\[
n_X=\frac{3\zeta(3)}{\pi^2}
\left[\frac{3\,\mathrm{MeV}}{\hbar c}\right]^3
\simeq1.29\times10^{39}\,\mathrm{m^{-3}}.
\]
(b) 由 $H=1/(2t)$ 及 Friedmann 式，$\rho=3/(32\pi Gt^2)$，在题目时刻约 $4.47\times10^9\,\mathrm{g\,cm^{-3}}$。(c) 能量自由度 $g_E=2+7/2+21/4+7/2=57/4$，由 $\rho c^2=g_E\pi^2(\kb T)^4/[30(\hbar c)^3]$ 得
\[
\kb T=\left[\frac{30\rho\hbar^3c^5}{\pi^2(57/4)}\right]^{1/4}
\simeq8.02\,\mathrm{MeV}.
\]
(d) 同一温度的额外辐射使 $H$ 较大，弱反应较早冻结，$n/p$ 较高；到氘瓶颈解除所用时间也变短，自由中子衰变较少，故理想模型的原初 $^4\mathrm{He}$ 增加。确切丰度需反应网络。(e) 解耦后的电磁浴在新粒子与电子对湮灭前有熵自由度 $2+7/2+7/2=9$，之后只有光子 $2$，故 $T_\nu/T_\gamma=(2/9)^{1/3}$。不能把中微子已解耦后的自身熵计入电磁浴。
\end{solution}

\src{8.286 Quiz 3 Review Problem 6(a)--(d), 原题为2009年 Quiz 3 Q3；官方复习解答。}
\begin{exercise}
假想弱作用理论 NTWI 使中微子在 $\kb T\simeq200\,\mathrm{MeV}$ 解耦，另有质量 $50\,\mathrm{MeV}/c^2$、各自旋 $1/2$ 的带电 $R^\pm$ 对。该温度的浴中还有光子、电子对、三味中微子、缪子对和三个自旋零的 $\pi^\pm,\pi^0$。忽略阈值的连续变化并假定空间平直：(a) 求标准模型和 NTWI 在该温度的质量密度；(b) NTWI 今天的 $T_\nu/T_\gamma$；(c) 两种理论今天三味热中微子加反中微子与光子数密度之比；(d) 假设弱反应在 $200\,\mathrm{MeV}$ 同时冻结，中子自由寿命约 $886\,\mathrm s$，估计氘瓶颈在通常约 $200\,\mathrm s$ 后打开时的氦质量份额方向与量级。给 $m_nc^2-m_pc^2=1.29\,\mathrm{MeV}$。
\end{exercise}
\begin{solution}
(a) 标准模型按题目指定的粒子计 $g_E=2+7/2+21/4+7/2+3=69/4$；NTWI 再加 $R^\pm$ 的 $7/2$，得 $83/4$。两者分别为
\[
\rho(g)=\frac{g\pi^2(200\,\mathrm{MeV})^4}{30\hbar^3c^5};
\quad \rho_{\rm std}\simeq2.10\times10^{18},\quad
\rho_{\rm NTWI}\simeq2.53\times10^{18}\ \mathrm{kg\,m^{-3}}.
\]
能量单位转换成焦耳后才可用 SI 的 $\hbar,c$。把所有粒子当严格无质量是题目的阶跃近似，真实 $200\,\mathrm{MeV}$ 附近不能忽略强相变细节。
(b) 解耦后仍与光子热耦合的浴有 $g_{s,\mathrm{before}}=2+7/2+7/2+3+7/2=31/2$，最终为 $2$。所以 $T_\nu/T_\gamma=(4/31)^{1/3}$，标准结果为 $(4/11)^{1/3}$。
(c) $n_{\nu+\bar\nu}/n_\gamma=(9/4)(T_\nu/T_\gamma)^3$，标准是 $9/11$，NTWI 是 $9/31$。
(d) 初始 $n/p=\exp(-1.29/200)\simeq0.994$，接近 1，故比标准模型留下更多中子。将中子分数近似取 $1/2$，在氘瓶颈之前仅考虑自由衰变，则 $Y_4\simeq\exp(-t_{\rm nuc}/886\,\mathrm s)$。在固定瓶颈温度下，后期辐射自由度 $g_{\rm eff}=2+\frac{21}4(T_\nu/T_\gamma)^4$；NTWI 的 $g_{\rm eff}$ 较小，故 $t_{\rm nuc,NTWI}\simeq200\sqrt{g_{\rm eff,std}/g_{\rm eff,NTWI}}\simeq240\,\mathrm s$，得到 $Y_4\simeq0.76$。它仅是题目刻意简化的数量级结果，不应当作真实可行弱理论的精确 BBN 预测。
\end{solution}

\src{8.286 Quiz 3 Review Problem 7, 原题为2011年 Quiz 3 Q4；官方复习解答。}
\begin{exercise}
除标准电子与正电子外，若再有一对质量与性质相同的电子样粒子及反粒子，假设中微子先解耦，今天中微子与 CMB 光子的温度比是多少？
\end{exercise}
\begin{solution}
两套带电费米子各贡献电磁浴熵自由度 $7/2$，光子贡献 $2$，湮灭前共 $9$，湮灭后为 $2$。两浴各自共动熵守恒，给
$T_\nu/T_\gamma=(2/9)^{1/3}$。
数值虽与上一题额外 $0.750\,\mathrm{MeV}$ 粒子同式，原因是两题对\emph{解耦后电磁浴的有效自由度}设定相同，并非质量细节普遍无关。
\end{solution}
